% Short running form: the full title is wider than the box the running head
% leaves beside the Sunday's name.
\runninghead{Render to God the One Who Bears His Image}
\properlane{image-and-inscription}{Render to God the One Who Bears His Image: The Coin, Its Image and the Self Owed to God}
\label{sec:image-and-inscription}

Read from its last sentence, the Gospel asks a question it does not answer
aloud. Jesus makes his questioners show the coin and name the image on it, and
from their answer he concludes that the coin is Caesar's. He then says that
God too is to be given what is his, and stops. This interpretation follows the
inference several Fathers draw from that silence: as the coin bears Caesar's
image and so belongs to Caesar, the human person bears God's image and so
belongs to God. Augustine draws the inference from the coin to the soul made
in God's image; Ambrose, on Luke's form of the saying, says that Christ bears
no image of Caesar because he is the image of God; Hilary of Poitiers argues
from a different premise, that body, soul and will are owed to God because
they come from him; Aquinas keeps a literal reading and gives the image reading
as the mystical sense. Carried through the whole Mass, the inference makes
``the things that are God's'' not a second tax set beside Caesar's but the
self, and the Mass's prayers ask God to receive it as a devoted will and an
undivided heart.

\subsection{What the coin settles, and what it leaves open}

The questioners asked about one debt, the census owed to Caesar, and the
answer speaks of two. The first is settled by the coin itself. Jesus does not
tell them whose the denarius is; he asks, and they say. The coin carries an
image and an inscription, and both name its owner. Whatever Jesus may think of
the census, the questioners already carry Caesar's coin and use it, and the
coin declares whose it is.

The second debt is not settled by anything in the scene. ``To God, the things
that are God's'' names no object. A reader who wants to know what is owed to
God must supply it, and here is the difficulty on which this interpretation
turns: Jesus never says ``the image of God''. Nothing in Matthew's words moves
from the face on the coin to the human face. The step is an inference, and it
depends on another text, the account of man's creation: ``Let us make man to
our image and likeness'' (Gn 1:26). The question for this interpretation is
whether that inference is warranted, and what follows if it is.

\subsection{Augustine: God looks for his image in man}

Augustine takes the step in at least five places and states its logic each
time. In his exposition of Psalm~4 he introduces it carefully. The psalm says
that the light of God's countenance is ``stamped on us'', and Augustine
explains, ``as a penny is stamped with the king's image. For man was made after
the image and likeness of God, which he defaced by sin: therefore it is his true
and eternal good, if by a new birth he be stamped.'' Then he adds: ``And I
believe this to be the bearing of that which some understand skilfully; I
mean, what the Lord said on seeing Caesar's tribute money \ldots\ As if He had
said, In like manner as Caesar exacts from you the impression of his image, so
also does God: that as the tribute money is rendered to him, so should the
soul to God, illumined and stamped with the light of His countenance''
(\work{Enarrationes in Psalmos} 4.8). Three things are visible in that passage.
The inference is not presented as his discovery; ``some'' understand it so,
and he endorses them. It rests on Genesis, on man's making after God's image.
And it does not stop at the image as given: the image was defaced by sin, and
the good of man is to be stamped again ``by a new birth''.

In his tractates on John the same inference becomes a sentence: ``For as Caesar
looks for his own image on the coin, so God looks for His in man''
(\work{In Iohannis evangelium tractatus} 41.2). In the tractate before it he
develops the defacement into a parable of his own. ``We are God's money: we
have wandered away as coin from the treasury. The impression that was stamped
upon us has been rubbed out by our wandering. He has come to refashion, for He
it was that fashioned us at first; and He is Himself asking for His money, as
Caesar for his. Therefore He says, `Render unto Caesar the things that are
Caesar's, and unto God the things that are God's:' to Caesar his money, to God
yourselves. And then shall the truth be reproduced in us'' (40.9). Here the
coin has a history. It was minted in God's treasury, it strayed, its stamp wore
smooth, and the one who struck it has come to strike it again. What God asks
is the coin itself, ``yourselves'', and he asks it as the one who can restore
what he asks for.

Twice more Augustine turns the inference into a command. In the exposition of
Psalm~57: ``Caesar seeketh his image; render it: God seeketh His image; render
it. Let not Caesar lose from you his coin: let not God lose in you His coin''
(\work{Enarrationes in Psalmos} 57.8). And on the psalm verse ``I will pay my
vows to the Lord'', he asks what vows a man should pay, and answers: ``let him
vow himself, let him pay himself as a vow: this is exacted, this is due. On
looking at the coin, the Lord saith, `Render unto Caesar the things which are
Caesar's, and unto God the things which are God's:' his own image is rendered
unto Caesar: let His image be rendered unto God'' (115.13). The person is not
asked to give God something of his own; he is asked to give himself back, as
a debt.

\subsection{Ambrose: the image Christ does not bear}

Ambrose reads the saying in Luke's Gospel, where it stands in the same
exchange, and gives it a centre that Augustine's passages imply but do not
state. Asked for the coin, he says, the Lord asks about the image, ``for the
image of God is one thing, the image of the world another'' (\latin{alia enim
imago Dei, alia imago mundi}), and he cites Paul: ``as we have borne the image
of the earthly, let us bear also the image of the heavenly'' (1~Cor 15:49).
Then he applies it to Christ himself: \latin{Imaginem Caesaris non habet
Christus, quia imago est Dei}, ``Christ does not have Caesar's image, because
he is the image of God''; nor does Peter, who left all to follow him
(\work{Expositio evangelii secundum Lucam} IX.34--35). Why, then, did Christ
pay the census? \latin{Non de suo dedit, sed reddidit mundo quod erat mundi}:
``he did not give of his own, but returned to the world what was the
world's''.

Ambrose's argument moves the image inference from the believer to its source.
Christ is the image of God in person, and so he owes Caesar nothing of his
own; the believer bears the image of the heavenly to the extent that he has
left what is the world's. Ambrose ends the passage with the believer's
renunciation. ``It is well that he decrees first that the things that are
Caesar's be rendered to Caesar: for no one can be the Lord's unless he has
first renounced the world'' (\latin{neque enim potest esse quis Domini, nisi
prius renuntiaverit mundo}). All renounce it in words when they receive the
sacraments, he says, but not all in affection, and then: \latin{Maior est
contractus fidei, quam pecuniae}, ``the contract of faith is greater than the
contract of money'' (IX.36). What is owed to God is a promise made at the
font, to be paid ``while you are in this body'' (\latin{Redde promissum, dum
in hoc corpore es}).

\subsection{Hilary: owed to the one from whom it comes}

Hilary of Poitiers, commenting on Matthew, reaches the same conclusion by
another road. He marvels at the answer: \latin{O plenam miraculi
responsionem}, ``O answer full of wonder''. It so balances the contempt of the
world against any offence to Caesar that it frees minds devoted to God
(\latin{devotas Deo mentes}) from all human cares while decreeing that Caesar
be given what is his. And what is God's own? \latin{Deo autem quae eius sunt
propria reddere nos oportere, corpus, animam, voluntatem}: we must render to
God what is properly his, ``body, soul, will''. His reason is not the image:
\latin{Ab eo enim haec profecta atque aucta retinemus: et proinde condignum
est, ut ei se totum reddant, cui debere se recolunt et originem et profectum},
``for we hold these as having come from him and been increased by him; and so
it is fitting that men render themselves whole to him to whom they recall that
they owe both their origin and their growth'' (\work{Commentarius in
Matthaeum} 23.2, PL~9, col.~1045).

Hilary's premise is origin and increase; Augustine's is the image stamped and
restamped. The two arguments arrive at the same object, the whole self, from
different starting points. For Augustine the coin's image is the hinge; for
Hilary the coin is Caesar's because it is in use among those who handle
Caesar's goods, and the self is God's because it came from God. His list also
names something Augustine's sentences do not: the will.

\subsection{Aquinas: the literal and the mystical sense}

Aquinas, commenting on Matthew, gives the saying two senses and holds them
together. Literally, he says, men have natural riches from God, such as bread
and wine, and of these they should give to God; they have artificial things,
such as denarii, from Caesar, and these they should render to Caesar
(\latin{Habetis divitias naturales a Deo, scilicet panem, et vinum; et de his
date Deo. Habetis ista artificialia, ut denarios, a Caesare; et haec Caesari
reddite}). Then: \latin{Mystice sic. Nos habemus animam quae est ad imaginem
Dei, ideo eam Deo reddere debemus}, ``Mystically thus: we have a soul which is
to the image of God, and therefore we owe it to God'' (\work{Super Matthaeum}
c.~22, Venice ed., p.~284). In Aquinas the two readings do not compete. The
literal sense gives God natural goods; the mystical sense gives him the soul
that bears his image.

\subsection{The strongest alternative: Jerome's tithes and offerings}

Jerome reads the second half of the saying literally. ``Render to Caesar the
things that are Caesar's, that is, coin, tribute and money; and to God the things that are
God's: let us understand tithes, first-fruits, oblations and victims''
(\latin{Reddite quae sunt Caesaris Caesari, id est, nummum tributum, et
pecuniam: et quae sunt Dei Deo, decimas primitias, et oblationes, ac victimas
sentiamus}). He adds that Christ himself did both: he paid the tribute for
himself and for Peter, and he rendered to God the things that are God's
``doing the will of the Father'' (\latin{Patris faciens voluntatem};
\work{Commentarii in Matthaeum} III, PL~26, col.~163).

On Jerome's reading the image belongs to the coin alone. The saying sets the
dues of the temple beside the dues of the empire, and it is a ruling on two
obligations, each with its own object. That is the plain sense of the words.
The image reading is the Fathers' spiritual reading of the coin, resting on
Genesis and not on anything Jesus says, and Augustine himself credits it to
``some''. Aquinas, as above, keeps a literal sense of his own beneath it.

Jerome's last clause, however, comes close to the others by another route.
Christ rendered to God ``doing the will of the Father''; Hilary names the will
among the things owed to God. On that point the literal reader and the reader
of origins meet without the image: what God is owed from Christ, and from those
who follow him, includes the will. Cyril of Alexandria, preaching on Luke's
form of the saying, says something similar: those who govern impose a tribute of money on their subjects, ``but God
requires of us of things corruptible and transitory even nothing, but rather
willing obedience and submission; faith and love; and the sweet savour of good
works'' (\work{Commentary on Luke}, Sermon 135).

\subsection{The Missal's texts in this reading}

Read beside the Gospel, the Collect asks for exactly the will that Hilary names
among the things owed to God, and that Jerome says Christ rendered in doing the
Father's will. It asks that the assembly always hold its will devoted to God
and serve his majesty \latin{sincero corde}. On this reading the Gospel's last
words become the measure of the Collect's petition.

The Entrance Antiphon gives the self that is rendered the posture of trust. The
one who prays asks to be kept ``as the apple of thy eye'', the small part by
which, in Augustine's explanation, the eye tells light from darkness
(\work{Enarrationes in Psalmos} 16.8); Aquinas reports the Gloss that reads
the pupil as Christ, who discerns the faithful from the unfaithful. On this
reading the person who is God's coin asks to be guarded as God guards what is
most his own.

The first reading contributes least. In the Clementine Vulgate the Lord says
to Cyrus, \latin{vocavi te nomine tuo: assimilavi te}, ``I have called thee by
thy name: I have made a likeness of thee''. Jerome comments that Cyrus was
called by name lest he be thought to be the Christ himself, ``to whom you were
likened, and of whom you went before as a type and image'' (\latin{cui
assimilatus es, et in cuius typum et imaginem praecessisti}; \work{Commentarii
in Isaiam} XII, PL~24, col.~457). A likeness given by God and a name given by
God would frame the Gospel's image and inscription. But the frame rests on the
Vulgate's word and on Jerome alone. The Lectionary's English, following the
Hebrew, does not speak of a likeness, and Aquinas reads \latin{assimilavi} of
the other saints whose names were foretold, without a typology.

The psalm says what the self owes when it enters God's courts. ``Bring ye to
the Lord, O ye kindreds of the Gentiles, bring ye to the Lord glory and
honour'' (Ps 95:7). Augustine says: ``Ascribe them not unto yourselves:
because they also who have declared it unto you, have not declared their own,
but His honour.'' On the courts of v.~8: ``Confession is a present unto God. O
heathen, if ye will enter into His courts, enter not empty \ldots\ Enter with
an humble heart into the house of God, and thou hast entered with a present''
(\work{Enarrationes in Psalmos} 95.8--9). The gift the nations bring is
themselves, humbled.

The second reading describes what belonging to God looks like when it is
lived. The church of the Thessalonians is ``in God the Father and in the Lord
Jesus Christ'', a phrase Chrysostom takes as a description of how one walks
(\work{Homilies on First Thessalonians} 1). In this reading the church's work
of faith, labour of charity and enduring hope are the self rendered in deeds.

The Alleluia verse adds a second figure of borrowed likeness. The faithful
shine as lights, and Aquinas says that they are light only by having something
of the light of the one who was the light of men (\work{Super ad Philippenses}
c.~2, lect.~4). Set beside the coin, that light behaves like an image: what
shines in them is another's.

The Prayer over the Offerings calls the gifts on the altar God's own
(\latin{donis tuis}) before they are offered, and asks that they be served
with a free mind. Beside it, Jerome's oblations and Aquinas's bread and wine
given to God become the literal counterpart of this reading: what is rendered
at the altar was God's before it was offered, as the self is.

The first Communion antiphon says that the eyes of the Lord rest on those who
fear him. Augustine's God looks for his image in man (\work{Tractatus} 41.2);
the antiphon names the God who looks. Where the second antiphon is sung, it
names the one who rendered himself entirely: the Son of Man came to give his
life as a ransom. Ambrose's Christ owes Caesar nothing of his own because he
is the image of God; in this antiphon the same Christ gives his life. Bede
reads the verse of the one who took the form of a servant to pour out his
blood for the world.

The Prayer after Communion asks that the assembly be aided by temporal
benefits and instructed by eternal ones. Read in this interpretation, the coin
belongs to the first kind and the self to the second: temporal goods help, but
the person is formed by what is eternal. Aquinas's two senses fall on either
side of that line, natural goods given to God and the soul that bears his
image.

\subsection{Where this interpretation strains}

The interpretation is a spiritual reading, and its strain is the distance
between the saying and the inference. Jesus speaks of a coin and of God; he
does not speak of man as God's coin. The Fathers who draw the inference supply
Genesis, and Augustine supplies a history of defacement and restamping that
Matthew does not narrate. A reader persuaded by Jerome will say that the
saying concerns two sets of dues and nothing more. The answer is not to deny
the literal sense, which Aquinas holds beside the mystical one, but to observe
that the literal sense leaves ``the things that are God's'' open, and that the
Fathers who fill the phrase with the self have a reason in Scripture for doing
so.

\subsection{The four senses of this interpretation}

\begin{fourSenses}
\item[Literal.] Asked whether tribute is lawful, Jesus makes his questioners
show the census coin and name the image and inscription it bears, and commands
that Caesar be repaid what is Caesar's and God what is God's (Mt 22:19--21).
On the literal reading the Fathers differ about the second object: tithes,
first-fruits, oblations and victims (Jerome), or body, soul and will
(Hilary).
\item[Allegorical.] Christ is the image of God, who bears no image of Caesar
(Ambrose), come to restamp the coin that wandering had worn smooth (Augustine,
\work{Tractatus} 40.9); where the second Communion antiphon is sung, he is the
Son of Man who gives his life as a ransom.
\item[Moral.] Render yourself: not a portion beside the tax, but the will
devoted and the heart undivided (the Collect), the body, soul and will owed to
their origin (Hilary), the person paying himself as a vow (Augustine on
Psalm 115), entering God's courts with a humble heart (Augustine on
Psalm 95).
\item[Anagogical.] The image restored is seen at the end. Augustine joins the
saying to the promise that ``we shall be like Him \ldots\ for we shall see Him
as He is'' (1~Jn 3:2) and says that ``then shall the truth be reproduced in
us'' (\work{Tractatus} 40.9); the soul is rendered to God ``stamped with the
light of His countenance'' (\work{Enarrationes in Psalmos} 4.8).
\end{fourSenses}
